\documentclass[UTF8,11pt]{ctexart}
\usepackage[a4paper,margin=24mm]{geometry}
\usepackage{amsmath,amssymb,amsthm,mathtools}
\usepackage[all]{xy}
\usepackage{xurl,hyperref,enumitem}
\hypersetup{hidelinks,breaklinks=true}
\setlength{\emergencystretch}{3em}
\newtheorem*{theorem}{Theorem}
\newtheorem*{lemma}{Lemma}
\newtheorem*{proposition}{Proposition}
\newtheorem*{corollary}{Corollary}
\theoremstyle{definition}
\newtheorem*{definition}{Definition}
\newtheorem*{remark}{Remark}
\DeclareMathOperator{\Hom}{Hom}
\DeclareMathOperator{\Ext}{Ext}
\DeclareMathOperator{\Tor}{Tor}
\DeclareMathOperator{\Spec}{Spec}
\DeclareMathOperator{\supp}{supp}
\DeclareMathOperator{\im}{im}
\DeclareMathOperator{\id}{id}
\DeclareMathOperator{\Tr}{Tr}
\DeclareMathOperator{\rank}{rank}
\newcommand{\R}{\mathbb R}
\newcommand{\C}{\mathbb C}
\newcommand{\Z}{\mathbb Z}
\newcommand{\N}{\mathbb N}
\newcommand{\Q}{\mathbb Q}
\begin{document}
\section*{012\quad 三维整数格暂留性的有限能量流证明}
\subsection*{来源与收录范围}
Peter G. Doyle and J. Laurie Snell, \emph{Random Walks and Electric Networks}, 5 July 2006. 主目标取自 \S2.4.7 标题及正文：三维整数格上的简单随机游走是暂留的。收录该节三维流的构造与能量估计（印刷 p.~114），连同 \S2.4.6 的能量上界定理和证明（p.~111）。后者只作支撑，不另计一道题。
\par\url{https://math.dartmouth.edu/~doyle/docs/walks/walks.pdf}
\par 获取日期：2026-09-28。已读取相关 PDF 文本并查看 pp.~110、114 的原页图像及 Figure~71。原始 PDF 未保存成功，附的是助手转录摘录。
\paragraph{记号与取材边界（整理者说明）}
三维格各边电阻为 $1$。原书 \S1.3.5（p.~49）定义流时采用 $j_{xy}=-j_{yx}$；无穷网络中除源点外满足 $\sum_yj_{xy}=0$，单位流满足 $\sum_yj_{0y}=1$，能量为 $\frac12\sum_{x,y}j_{xy}^2R_{xy}$。$G(r)$ 是源点周围的有限耗尽子图，边界接至共同终点；$R_{\mathrm{eff}}$ 是对应有效电阻的极限。这些记号据原书整理，不是新增证明。\S2.4.6 在 p.~110 重述定义时实际印成 $j_{xy}=j_{yx}$，与其 p.~49 的定义不同；本文说明采用前文的反对称约定，没有把改过的式子伪装成原句。
\subsection*{作者命题与人类证明}
\begin{theorem}[Section 2.4.6]
The effective resistance $R_{\mathrm{eff}}$ from $0$ to $\infty$ is less than or equal to the energy dissipation of any unit flow from $0$ to infinity.
\end{theorem}
\begin{proof}
Assume that we have a unit flow $j$ from $0$ to infinity with energy dissipation
\[
E=\frac12\sum_{x,y}j_{xy}^2R_{xy}.
\]
We claim that $R_{\mathrm{eff}}\le E$. Restricting $j_{xy}$ to the edges of the finite graph $G(r)$, we have a unit flow from $0$ to $\partial G(r)$ in $G(r)$. Let $i^{(r)}$ be the unit current flow in $G(r)$ from $0$ to $\partial G(r)$. By the results of Section 1.3.5,
\[
R_{\mathrm{eff}}^{(r)}=\frac12\sum_{G(r)}(i^{(r)}_{xy})^2R_{xy}
\le\frac12\sum_{G(r)}j_{xy}^2R_{xy}
\le\frac12\sum_{x,y}j_{xy}^2R_{xy}=E,
\]
where $\sum_{G(r)}$ indicates the sum over all pairs $x,y$ such that $xy$ is an edge of $G(r)$.
\end{proof}
\paragraph{A proof, using flows, that simple random walk in three dimensions is transient (Section 2.4.7).}
We now apply this form of the cutting method to give another proof that simple random walk on the threedimensional lattice is transient. All we need is a flow to infinity with finite dissipation.
\begin{proof}[Three-dimensional construction and energy estimate]
For three dimensions, the uniform flow is defined as follows: Out of $(x,y,z)$ with $x+y+z=n$ we have the flow indicated in Figure 71. The total flow out of $(x,y,z)$ is then
\[
\frac{2(x+1)}{(n+3)(n+2)(n+1)}+
\frac{2(y+1)}{(n+3)(n+2)(n+1)}+
\frac{2(z+1)}{(n+3)(n+2)(n+1)}
=\frac2{(n+2)(n+1)}.
\]
\begin{center}
\begin{tabular}{c|c}
\multicolumn{2}{c}{Figure 71: edge labels transcribed from the original diagram}\\
Outgoing direction & Flow \\\hline
positive $x$ & $\displaystyle\frac{2(x+1)}{(n+3)(n+2)(n+1)}$\\[7pt]
positive $y$ & $\displaystyle\frac{2(y+1)}{(n+3)(n+2)(n+1)}$\\[7pt]
positive $z$ & $\displaystyle\frac{2(z+1)}{(n+3)(n+2)(n+1)}$
\end{tabular}
\end{center}
The flow into $(x,y,z)$ comes from the points $(x-1,y,z)$, $(x,y-1,z)$, $(x,y,z-1)$ and, hence, the total flow into $(x,y,z)$ is
\[
\frac{2x}{(n+2)(n+1)n}+
\frac{2y}{(n+2)(n+1)n}+
\frac{2z}{(n+2)(n+1)n}
=\frac2{(n+2)(n+1)}.
\]
Thus the net flow for $(x,y,z)$ is $0$. The flow out of $0$ is $(1/3)+(1/3)+(1/3)=1$ We have now to check finiteness of energy dissipation. The flows coming out of the edges at the $n$th level are all $\le2/(n+1)^2$. There are $(n+1)(n+2)/2$ points a distance $n$ from $0$, and thus there are $(3/2)(n+1)(n+2)\le3(n+1)^2$ edges coming out of the $n$th level. Thus the energy dissipation $E$ has
\[
E\le\sum_n3(n+1)^2\left(\frac2{(n+1)^2}\right)^2
=12\sum_n\frac1{(n+1)^2}<\infty,
\]
and the random walk is transient.
\end{proof}
\subsection*{整理说明（非作者正文）}
正文的流位于正卦限，原书在三维构造前用二维同类流作教学说明；二维示例不参与三维能量收敛的证明，未收作第二个问题。Figure~71 的三个正方向及对应流量全部保留，改为表格式转录而非冒称保留原始图像。

标准先修为有限网络的 Thomson 原理、耗尽定义的有效电阻及其与暂留性的对应、向更大网络增边的电阻单调性。本题承担的是明确构造单位流、验证守恒、按距离层计算边数并证明总能量有限；不能把“三维格暂留”本身作先修。此题只给三维暂留性，不冒充种子题“任意无限图暂留当且仅当存在有限能量流”的双向定理。
\subsection*{可修改的简短标注（非作者正文）}
$c$：三维整数格随机游走的暂留性

$a$：单位流；有限网络Thomson原理；有效电阻；正卦限分层构造；能量求和

\texttt{cross\_theory\_status = yes}。将随机游走的暂留目标转成电网络中有限能量单位流的构造；层上的组合计数给出能量收敛，并返回暂留结论。位置：\S2.4.6--2.4.7，尤其 p.~114。

全部标注均为非穷尽、可修改初稿；未标注不表示未使用。
\subsection*{许可与修改记录}
原数学正文作者为 Peter G. Doyle 与 J. Laurie Snell；原书 2006 年版本按 GNU Free Documentation License 发布，无不变章节、封面文字或封底文字。许可全文随包保留在 \texttt{sources/stacks/GFDL-1.2.txt}。本题标题不同于原书，整理部分随原文许可；不暗示原作者认可本次标注。2026-09-28：助手转录、加入来源定位和说明；未新增数学证明。
\end{document}
